\documentclass[11pt]{article}
\usepackage[margin=1in]{geometry}
\usepackage[T1]{fontenc}
\usepackage[utf8]{inputenc}
\usepackage{array}
\usepackage{booktabs}
\usepackage{longtable}
\usepackage{hyperref}
\title{graphicsTools.wl Module Documentation}
\author{MFGraphs maintainers}
\date{}
\begin{document}
\maketitle

\section{High-level explanation}

\texttt{graphicsTools.wl} is the visualization layer of the active MFGraphs package. It exposes two plotters and one structural helper. \texttt{rawNetworkPlot} renders the physical network as the user sees it (real vertices and physical edges), and \texttt{richNetworkPlot} renders the \emph{augmented state-space graph} (oriented-edge nodes, flow arcs, and transition arcs) that the system builders actually solve over. \texttt{augmentAuxiliaryGraph} produces the underlying augmented-graph data that \texttt{richNetworkPlot} consumes, so external tooling can build its own views without re-deriving the augmented structure.

The module deliberately ships only those two visual idioms. Every styling concern --- whether boundary data is shown, which labels appear, how arcs are coloured, how anti-parallel edges are separated --- is exposed through plot options rather than through additional public functions.

\section{Public symbols}

\subsection*{\texttt{rawNetworkPlot}}

\texttt{rawNetworkPlot[s, sys, opts]} and \texttt{rawNetworkPlot[s, sys, sol, opts]} render the physical network. Real network vertices are drawn in gray; auxiliary entry/exit vertices and boundary edges are hidden by default. When a solution is supplied, shown vertex states are coloured by a $u$-value gradient (blue $\to$ red) while physical entry/exit vertices remain gray. Density labels can be overlaid for Tawaf-with-density systems.

Options (all defaults shown):

\begin{center}
\renewcommand{\arraystretch}{1.2}
\begin{longtable}{@{}p{4.2cm} p{2.4cm} p{8.4cm}@{}}
\toprule
\textbf{Option} & \textbf{Default} & \textbf{Effect} \\
\midrule
\endhead
\texttt{ShowAuxiliaryVertices} & \texttt{False} & Show the auxiliary entry/exit vertices added by \texttt{makeScenario}. \\
\texttt{ShowBoundaryData}      & \texttt{False} & Overlay boundary entry/exit data; forces \texttt{ShowAuxiliaryVertices} on. \\
\texttt{ShowBoundaryValues}    & \texttt{True}  & Show $u$-values on auxiliary exit vertices when boundary data is shown. \\
\texttt{ShowFlowLabels}        & \texttt{Automatic} & Label edges with $j$-values; \texttt{True} when a solution is supplied. \\
\texttt{ShowValueLabels}       & \texttt{Automatic} & Label vertex states with $u$-values; \texttt{True} when a solution is supplied. \\
\texttt{ShowDensityLabels}     & \texttt{False} & Show inferred per-edge densities $m$ (Tawaf-with-density). \\
\texttt{ColorFunction}         & \texttt{Automatic} & Custom colour function over $u$-values (default: blue-to-red blend). \\
\texttt{ShowLegend}            & \texttt{True}  & Display the legend. \\
\texttt{GraphLayout}           & \texttt{Automatic} & Forwarded to the underlying \texttt{Graph} call. \\
\texttt{PlotLabel}             & \texttt{Automatic} & Title styling. \\
\texttt{ImageSize}             & \texttt{Large} & Forwarded to the underlying \texttt{Graph} call. \\
\bottomrule
\end{longtable}
\end{center}

\subsection*{\texttt{richNetworkPlot}}

\texttt{richNetworkPlot[s, sys, opts]} and \texttt{richNetworkPlot[s, sys, sol, opts]} render the augmented state-space graph. Nodes are oriented-edge pairs $\{a, b\}$ representing the state ``arriving at $b$ along edge $(a, b)$''. Edges come in two colours: flow arcs $j[a, b]$ in blue and transition arcs $j[r, i, w]$ in red. Arcs are drawn as quadratic Bezier curves so that anti-parallel pairs separate visually. Only nodes whose second component is an auxiliary entry/exit vertex receive boundary colouring.

Options (all defaults shown):

\begin{center}
\renewcommand{\arraystretch}{1.2}
\begin{longtable}{@{}p{4.2cm} p{2.4cm} p{8.4cm}@{}}
\toprule
\textbf{Option} & \textbf{Default} & \textbf{Effect} \\
\midrule
\endhead
\texttt{ShowFlowEdges}      & \texttt{True}  & Include flow arcs $j[a, b]$. Set to \texttt{False} to obtain the transition-only graph. \\
\texttt{ShowBoundaryData}   & \texttt{False} & Overlay entry-flow and exit-cost labels. \\
\texttt{ShowBoundaryValues} & \texttt{True}  & Show $u$ / exit-cost on boundary nodes. \\
\texttt{ShowFlowLabels}     & \texttt{Automatic} & Label arcs with $j$-values; \texttt{True} when a solution is supplied. \\
\texttt{ShowValueLabels}    & \texttt{Automatic} & Label nodes with $u$-values; \texttt{True} when a solution is supplied. \\
\texttt{UseColorFunction}   & \texttt{False} & Colour nodes and eligible edges by $u$-values via \texttt{ColorFunction}. \\
\texttt{ColorFunction}      & \texttt{Automatic} & Custom colour function over $u$-values. \\
\texttt{ShowLegend}         & \texttt{True}  & Display the legend. \\
\texttt{BendFactor}         & \texttt{0.15}  & Arc curvature as a fraction of edge length. \\
\texttt{GraphLayout}        & \texttt{Automatic} & Forwarded to the underlying \texttt{Graph} call. \\
\texttt{PlotLabel}          & \texttt{Automatic} & Title styling. \\
\texttt{ImageSize}          & \texttt{Large} & Forwarded to the underlying \texttt{Graph} call. \\
\bottomrule
\end{longtable}
\end{center}

\subsection*{\texttt{augmentAuxiliaryGraph}}

\texttt{augmentAuxiliaryGraph[sys]} constructs the augmented state-space graph from the system's \texttt{AuxPairs} and \texttt{AuxTriples}. It returns an \texttt{Association} with keys:
\begin{itemize}
\item \texttt{"Graph"} --- the Wolfram \texttt{Graph} object;
\item \texttt{"Vertices"} --- the list of oriented-edge node labels;
\item \texttt{"FlowEdges"} --- flow arcs corresponding to the $j[a, b]$ family;
\item \texttt{"TransitionEdges"} --- transition arcs corresponding to the $j[r, i, w]$ family;
\item \texttt{"EdgeVariables"} --- map from each arc to the symbolic variable it represents;
\item \texttt{"EdgeKinds"} --- map from each arc to either \texttt{"Flow"} or \texttt{"Transition"}.
\end{itemize}
This is the data \texttt{richNetworkPlot} renders. External tooling that wants to draw its own augmented-state view (e.g.\ in a notebook, in a publication figure, or in an animation) should consume this association directly rather than reimplementing the augmentation logic.

\section{Usage guidance}

A typical plotting workflow is:

\begin{verbatim}
s   = gridScenario[{3}, {{1, 10}}, {{3, 0}}];
sys = makeSystem[s];
sol = solveScenario[s];

raw  = rawNetworkPlot[s, sys, sol];                  (* physical view *)
rich = richNetworkPlot[s, sys, sol];                 (* augmented view *)
transOnly = richNetworkPlot[s, sys, sol, ShowFlowEdges -> False];
\end{verbatim}

Choose \texttt{rawNetworkPlot} when the audience cares about the network as a road map or topology, and \texttt{richNetworkPlot} when they need to reason about transition structure and switching choices. The two plots are complementary: showing them side by side is often the clearest way to communicate what an MFGraphs solution means.

\section*{References}

\begin{enumerate}
\item Source file: \texttt{MFGraphs/graphicsTools.wl}
\item Upstream dependencies: \texttt{scenarioTools.wl}, \texttt{systemTools.wl}
\item Test suite: \texttt{MFGraphs/Tests/graphicsTools.mt}
\end{enumerate}

\end{document}
